\documentclass[10pt,a4paper]{amsart}
\usepackage[margin=19mm]{geometry}
\usepackage{amsmath,amssymb,mathtools}
\usepackage[scr=rsfs,cal=euler]{mathalfa}
\usepackage{xfrac}
\usepackage[unicode,hidelinks,bookmarks=false]{hyperref}
\newcommand{\ie}{i.e.\ }
\newcommand{\cf}{cf.\ }
\newcommand{\resp}{respectively\ }
\newcommand{\Subsection}{\subsection}
\providecommand{\nobreakdash}{\nobreak\hbox{-}}
\setlength{\emergencystretch}{2em}
\multlinegap0pt
\newcommand{\e}{\mathrm{e}}
\theoremstyle{plain}
\newtheorem{theorem}{Theorem}
\newtheorem{proposition}{Proposition}
\newtheorem{conjecture}{Conjecture}
\newtheorem{problem}{Problem}
\theoremstyle{remark}
\newtheorem{remark}{Remark}
\title[Morse Topology, Zero Products, and Elliptic Prime Crossings]{Morse Topology, Zero Products, and Elliptic Prime Crossings}
\author{Michel Planat}
\date{Source: arXiv:2609.27962v1 (CC BY 4.0)}
\begin{document}
\maketitle

\noindent\emph{Source and licence.} The delivered work is the article shown in the title above, version arXiv:2609.27962v1 (CC BY 4.0), distributed under the Creative Commons Attribution 4.0 International licence, as recorded in the arXiv licence metadata for this exact version and shown on the abstract page saved with this item. The full licence notice, which carries the licence URL and the address of that abstract page, is the bundled file \texttt{licenses/arxiv-2609.27962-CC-BY-4.0.txt}.
\par\smallskip

\noindent\textit{Editorial scope.} Counted unit: the article's Theorem (source label eq:hard-edge-tail-law) (source0.tex lines 1105--1125, source label \texttt{eq:hard-edge-tail-law}): ``Odd-orthogonal hard-edge tail For every fixed , there is a finite constant such that J N y C N -3y 2 , y . More explicitly, with , C N N 48Z N (0, ) N-1 A N( u) ( 32 k 2 N j(u k) )d u, A N( u) k 2 N (''. It is delivered together with the article's complete original proof (source0.tex lines 1127--1155, one proof environment adjacent to the statement, 29 lines), copied line for line from the article's own text. Required context retained uncounted: eq:SOodd-Weyl-density (lines 1091--1097). Editorial formatting only: an AMS \texttt{amsart} preamble that restates the article's own macro definitions and its own theorem declarations, and the theorem environments the retained lines use; the article's own class file is neither shipped nor input; the retained environments are renumbered sequentially in order of appearance, whereas the article prints them with its own numbering; the article's equation numbering is not reproduced and every cross-reference inside the retained text resolves to the wrapper's numbering. No citation occurs inside any retained line, so the wrapper carries no bibliography; All retained bodies are exact copies of the bundled \texttt{sources/211535/source0.tex}, so no omission receipt applies. No asset-inclusion command occurs inside any retained span and the package ships no image asset.


\section*{Retained context: eq:SOodd-Weyl-density (source0.tex lines 1091--1097)}

\begin{equation}
  p_N(\boldsymbol\theta)
  =\frac1{Z_N}\prod_{j=1}^{N}(1-\cos\theta_j)
   \prod_{1\le j<k\le N}(\cos\theta_j-\cos\theta_k)^2,
  \qquad \boldsymbol\theta\in(0,\pi)^N,
  \label{eq:SOodd-Weyl-density}
\end{equation}


\section{The counted unit: Theorem (source label eq:hard-edge-tail-law) and its complete original proof}

\begin{theorem}[Odd-orthogonal hard-edge tail]
For every fixed $N\ge1$, there is a finite constant $C_N>0$ such that
\begin{equation}
  \Pr\{\mathcal J_N>y\}
  \sim C_N\e^{-3y/2},
  \qquad y\to\infty.
  \label{eq:hard-edge-tail-law}
\end{equation}
More explicitly, with $\mathbf u=(u_2,\ldots,u_N)$,
\begin{align}
  C_N&=\frac{N}{48Z_N}\int_{(0,\pi)^{N-1}}
  A_N(\mathbf u)
  \exp\!\left(\frac32\sum_{k=2}^{N}j(u_k)\right)d\mathbf u,
  \label{eq:hard-edge-constant}\\
  A_N(\mathbf u)&=
  \prod_{k=2}^{N}(1-\cos u_k)^3
  \prod_{2\le j<k\le N}(\cos u_j-\cos u_k)^2.
  \label{eq:hard-edge-boundary-density}
\end{align}
For $N=1$, the integral is interpreted as $1$.
\end{theorem}

\begin{proof}
Distinguish one eigenangle $t$ and write the others as $\mathbf u$.  From \eqref{eq:SOodd-Weyl-density},
\begin{equation}
  p_N(t,\mathbf u)
  =\frac{t^2}{2Z_N}A_N(\mathbf u)(1+o(1)),
  \qquad t\downarrow0,
  \label{eq:Weyl-boundary-expansion}
\end{equation}
for almost every $\mathbf u$.  At the same time,
\begin{equation}
  j(t)=-2\log(2t)+o(1).
  \label{eq:j-hard-edge-asymptotic}
\end{equation}
Put $R(\mathbf u)=\sum_{k=2}^{N}j(u_k)$.  For fixed $\mathbf u$, the boundary of
$j(t)+R(\mathbf u)>y$ is
\begin{equation}
  t_y(\mathbf u)
  =\frac1{2\sqrt{\exp(y-R(\mathbf u))-1}}
  \sim\frac12\exp\!\left[-\frac{y-R(\mathbf u)}2\right].
\end{equation}
Integrating \eqref{eq:Weyl-boundary-expansion} from $0$ to $t_y(\mathbf u)$ gives
\begin{equation}
  \frac{A_N(\mathbf u)}{48Z_N}
  \exp\!\left(\frac32R(\mathbf u)\right)
  \e^{-3y/2}(1+o(1)).
\end{equation}
The boundary integral is finite.  Indeed, as one of the remaining angles tends to zero,
$(1-\cos u)^3\asymp u^6$ while $\exp(3j(u)/2)\asymp u^{-3}$, leaving an integrable factor $u^3$; eigenangle collisions only introduce additional zeros through the squared Vandermonde.  Dominated convergence therefore applies.  Summing over the $N$ possible distinguished eigenangles gives \eqref{eq:hard-edge-constant}.  Regions in which two or more angles approach the hard edge on the same exponential scale are of lower order because of the squared Vandermonde factor, so they do not alter the leading constant.
\end{proof}

\end{document}
